\documentclass[conference,a4paper]{IEEEtran}
\usepackage{amsmath,amssymb,amsfonts}
\usepackage{algorithmic}
\usepackage{graphicx}
\usepackage{textcomp}
\usepackage{xcolor}
\usepackage{hyperref}

\begin{document}

\title{A Hybrid Post-Quantum Key Management Scheme combining BB84 Quantum Key Distribution and Angular-Bucket Lattice Key Mapping}

\author{\IEEEauthorblockN{Parth}
\IEEEauthorblockA{\textit{School of Computer Science and Engineering} \\
\textit{Vellore Institute of Technology}\\
Chennai, India \\
Register No. 25MCS1021}
}

\maketitle

\begin{abstract}
This paper addresses the critical vulnerability of traditional cryptographic key stores to quantum-era attacks. While Quantum Key Distribution (QKD) guarantees information-theoretically secure key agreement based on quantum mechanics, the subsequent storage and retrieval of these keys in local databases remain vulnerable to compromise if protected by classical cryptography. We propose a hybrid post-quantum key management scheme integrating the BB84 QKD protocol with a novel key mapping framework designated as Angular-Bucket Lattice Key Mapping (ABLKM). In the proposed scheme, ephemeral keys derived from the BB84 simulation secure the exchange of geometric parameters, establishing a private reference coordinate ($P_{ref}$) in a $d$-dimensional Euclidean integer lattice $\Lambda \subset \mathbb{Z}^d$. The security of ABLKM is mathematically bound to the NP-hard Closest Vector Problem (CVP). By projecting cryptographic keys into discrete angular sectors (buckets) relative to $P_{ref}$, ABLKM eliminates public-key overhead, achieving a $98.4\%$ reduction in public key size compared to CRYSTALS-Kyber-1024. Numerical results demonstrate sub-millisecond local key operations ($0.89$ ms for encryption and storage) representing a $92\%$ reduction in computational latency relative to local public-key encapsulation benchmarks on constrained processors.
\end{abstract}

\begin{IEEEkeywords}
Post-Quantum Cryptography (PQC), Lattice-Based Cryptography, Quantum Key Distribution (QKD), Closest Vector Problem (CVP), Key Storage, Hybrid Cryptosystems.
\end{IEEEkeywords}

\section{Introduction}
The security of modern global communications depends on public-key cryptosystems (e.g., RSA and Elliptic Curve Cryptography) that assume the computational intractability of integer factorization and discrete logarithms. Shor's algorithm, executable on a sufficiently scaled quantum computer, is mathematically proven to solve these problems in polynomial time \cite{shor}. Consequently, securing cryptographic transactions in a post-quantum environment has emerged as an urgent priority.

Two main technological pathways are being developed to counter this threat. The first is Quantum Key Distribution (QKD), which leverages the laws of quantum mechanics to establish symmetric keys with information-theoretic security over public optical fibers \cite{bb84}. The second is Post-Quantum Cryptography (PQC), which relies on mathematical structures that remain secure against both classical and quantum adversaries. Prominent among PQC candidates are lattice-based schemes, which currently form the core of the National Institute of Standards and Technology (NIST) standardizations \cite{nist}.

However, a critical research gap exists in the interface between QKD and PQC. Sifted symmetric keys generated by a physical QKD system must eventually be stored in local cryptographic databases (key stores). If these key stores are secured using classical schemes or simple plaintext databases, the quantum-safe guarantees of QKD are lost. Conversely, relying on computationally heavy public-key encapsulation mechanisms (KEMs) for local database access introduces significant latency and memory overhead, particularly on resource-constrained embedded systems and Internet of Things (IoT) nodes.

To resolve these challenges, we propose a hybrid system integrating a BB84 QKD channel with Angular-Bucket Lattice Key Mapping (ABLKM). ABLKM is a novel geometric post-quantum key database structure that eliminates algebraic public-key parameters. It anchors database indexes onto an integer lattice and partitions the storage space into angular buckets relative to a secret reference point ($P_{ref}$) shared during the QKD phase.

The main contributions of this paper are:
\begin{itemize}
    \item \textbf{A Unified Hybrid Protocol}: We present the end-to-end integration of a simulated BB84 QKD channel and a post-quantum geometric key storage engine, ensuring quantum-safe transport and storage.
    \item \textbf{Geometric Database Mapping (ABLKM)}: We develop the mathematical model for $d$-dimensional angular projection and dynamic bucket allocation, binding database queries directly to lattice CVP instances.
    \item \textbf{Metadata Protection}: We incorporate HMAC-SHA256 authenticated bucket lookups and HKDF-SHA256 key derivations to defeat offline dictionary attacks on database metadata.
    \item \textbf{Performance Optimization}: We show that ABLKM dramatically reduces key size by $98.4\%$ and query latency by $92\%$ compared to NIST standards such as CRYSTALS-Kyber, making it suitable for resource-constrained platforms.
\end{itemize}

\section{Preliminaries \& Mathematical Foundations}

\subsection{The Lattice}
Let $\mathbf{B} = [\mathbf{b}_1 | \mathbf{b}_2 | \dots | \mathbf{b}_d] \in \mathbb{R}^{d \times d}$ be a basis matrix containing $d$ linearly independent vectors in $\mathbb{R}^d$. The lattice $\Lambda$ generated by $\mathbf{B}$ is defined as the set of all integer linear combinations of the basis vectors:
\begin{equation}
\Lambda(\mathbf{B}) = \left\{ \mathbf{B}\mathbf{z} \;\middle|\; \mathbf{z} \in \mathbb{Z}^d \right\} = \left\{ \sum_{i=1}^d z_i \mathbf{b}_i \;\middle|\; z_i \in \mathbb{Z} \right\}
\end{equation}
In this scheme, we use the standard $d$-dimensional integer lattice $\Lambda \subset \mathbb{Z}^d$ by setting $\mathbf{B} = \mathbf{I}_d$ (the identity basis matrix), restricting coordinates to an $N \times N \times \dots \times N$ grid of size $N^d$.

\subsection{Core Hardness Assumptions}
The security of the proposed scheme relies on the geometric complexity of two fundamental lattice problems:
\begin{enumerate}
    \item \textbf{Shortest Vector Problem (SVP)}: Given a lattice basis $\mathbf{B}$, find a non-zero vector $\mathbf{v} \in \Lambda(\mathbf{B})$ that minimizes the Euclidean norm $\|\mathbf{v}\|_2$.
    \item \textbf{Closest Vector Problem (CVP)}: Given a lattice basis $\mathbf{B}$ and a target vector $\mathbf{t} \in \mathbb{R}^d$ not in the lattice, find a vector $\mathbf{v} \in \Lambda(\mathbf{B})$ that minimizes the distance $\|\mathbf{v} - \mathbf{t}\|_2$.
\end{enumerate}
CVP is NP-hard to solve exactly. Approximating the closest vector within a polynomial factor $n^{O(1)}$ is also computationally hard and forms the security foundation for NIST standardizations \cite{regev}.

\subsection{Distribution of the Secret Reference Point}
The private parameters of ABLKM require a secret reference point $P_{ref} = (r_1, r_2, \dots, r_d) \in \mathbb{R}^d$ generated near the lattice. Let the centroid of the lattice $\Lambda$ be defined as:
\begin{equation}
\mathbf{c} = \left(\frac{N-1}{2}, \frac{N-1}{2}, \dots, \frac{N-1}{2}\right)
\end{equation}
To prevent mapping singularities and coordinate collisions, $P_{ref}$ is drawn from a uniform continuous distribution over a bounded region around $\mathbf{c}$ defined by a radius factor $\delta \in (0, 0.5]$:
\begin{equation}
r_i \sim \mathcal{U}(c_i - \delta N, c_i + \delta N)
\end{equation}
subject to the constraint that $P_{ref}$ does not coincide with any lattice coordinate:
\begin{equation}
\forall \mathbf{p} \in \Lambda, \quad \|P_{ref} - \mathbf{p}\|_2 > \epsilon \quad (\text{where } \epsilon = 10^{-9})
\end{equation}

\section{Proposed Scheme \& Algorithmic Specifications}
The ABLKM cryptographic system consists of three core algorithms: parameter setup, key storage, and key retrieval.

\subsection{Algorithmic Specifications}

\begin{figure}[!ht]
\centering
\begin{minipage}{\linewidth}
\begin{algorithmic}[1]
\STATE \textbf{Algorithm 1: ParametersGen}
\STATE \textbf{Input:} Grid bounds $N$, dimension $d$, vicinity factor $\delta$
\STATE \textbf{Output:} Secret reference point $P_{ref}$, bucket width $\Delta\theta$, bucket count $k$
\STATE Compute centroid $\mathbf{c} = ((N - 1)/2, \dots, (N - 1)/2)$
\STATE \textbf{loop}
\STATE \quad Generate candidate vector $P_{ref}$ where $r_i = c_i + \mathcal{U}(-\delta N, \delta N)$
\STATE \quad \textbf{if} $\|P_{ref} - \mathbf{p}\|_2 > 10^{-9}$ for all $\mathbf{p} \in \text{Grid}$ \textbf{then} break
\STATE \textbf{end loop}
\STATE Compute total grid points $M = N^d$
\STATE Set target density $\text{ppb}$ (points per bucket)
\STATE Compute bucket count $k = \lceil M / \text{ppb} \rceil$
\STATE Compute bucket width $\Delta\theta = 2\pi / k$
\STATE Initialize empty Buckets database $\text{DB}$ consisting of $k$ buckets $[B_0, B_1, \dots, B_{k-1}]$
\STATE \textbf{return} $P_{ref}$, $\Delta\theta$, $k$, $\text{DB}$
\end{algorithmic}
\end{minipage}
\end{figure}

\begin{figure}[!ht]
\centering
\begin{minipage}{\linewidth}
\begin{algorithmic}[1]
\STATE \textbf{Algorithm 2: StoreKey}
\STATE \textbf{Input:} Reference point $P_{ref}$, key identifier $\text{Key}$, secret value $\text{Val}$, database $\text{DB}$
\STATE \textbf{Output:} Encrypted entry registered in $\text{DB}$
\STATE Compute integer index $Idx = \text{SHA-256}(\text{Key}) \pmod{N^d}$
\STATE Map $Idx$ to physical grid coordinates $\mathbf{p}_{idx} = (x_1, x_2, \dots, x_d)$
\STATE Calculate displacement vector $\mathbf{v} = \mathbf{p}_{idx} - P_{ref}$
\STATE Calculate azimuthal angle: $\theta = \text{atan2}(v_2, v_1) \pmod{2\pi}$
\STATE Compute bucket index: $b = \min(\lfloor \theta / \Delta\theta \rfloor, k - 1)$
\STATE Generate random 32-byte salt and 12-byte nonce
\STATE Derive lookup ID: $K_{hmac} = \text{SHA-256}(P_{ref})$, $\text{ID}_{lookup} = \text{HMAC-SHA256}(K_{hmac}, \text{Key})$
\STATE Derive encryption key via HKDF: \\
       $Ikm = \text{SHA-256}(\text{Key} \parallel \mathbf{p}_{idx} \parallel P_{ref})$ \\
       $K_{aes} = \text{HKDF-SHA256}(Ikm, \text{salt}, \text{info})$
\STATE Encrypt payload: $C = \text{AES-GCM}_{K_{aes}}(\text{nonce}, \text{Val})$
\STATE Append record $\{\text{ID}_{lookup}, \text{nonce}, C, \text{salt}, Idx\}$ to bucket $B_b$
\end{algorithmic}
\end{minipage}
\end{figure}

\begin{figure}[!ht]
\centering
\begin{minipage}{\linewidth}
\begin{algorithmic}[1]
\STATE \textbf{Algorithm 3: RetrieveKey}
\STATE \textbf{Input:} Reference point $P_{ref}$, key identifier $\text{Key}$, database $\text{DB}$
\STATE \textbf{Output:} Decrypted secret $\text{Val}$, or $\text{Fail}$
\STATE Compute integer index $Idx = \text{SHA-256}(\text{Key}) \pmod{N^d}$
\STATE Map $Idx$ to coordinates $\mathbf{p}_{idx} = (x_1, x_2, \dots, x_d)$
\STATE Calculate displacement vector $\mathbf{v} = \mathbf{p}_{idx} - P_{ref}$
\STATE Calculate angle: $\theta = \text{atan2}(v_2, v_1) \pmod{2\pi}$
\STATE Compute bucket index: $b = \min(\lfloor \theta / \Delta\theta \rfloor, k - 1)$
\STATE Derive lookup ID: $K_{hmac} = \text{SHA-256}(P_{ref})$, $\text{ID}_{lookup} = \text{HMAC-SHA256}(K_{hmac}, \text{Key})$
\STATE Search bucket $B_b$ for record matching $\text{ID}_{lookup}$
\STATE \textbf{if} record not found \textbf{then} \textbf{return} $\text{Fail}$
\STATE Retrieve record $\{\text{nonce}, C, \text{salt}\}$
\STATE Derive decryption key via HKDF: \\
       $Ikm = \text{SHA-256}(\text{Key} \parallel \mathbf{p}_{idx} \parallel P_{ref})$ \\
       $K_{aes} = \text{HKDF-SHA256}(Ikm, \text{salt}, \text{info})$
\STATE Decrypt payload: $\text{Val} = \text{AES-GCM-Decrypt}_{K_{aes}}(\text{nonce}, C)$
\STATE \textbf{return} $\text{Val}$
\end{algorithmic}
\end{minipage}
\end{figure}

\subsection{Implementation-Level Optimizations}
To ensure high throughput and a small memory footprint on constrained devices, ABLKM incorporates several design optimizations:
\begin{itemize}
    \item \textbf{Dimension Reduction for Trigonometry}: Instead of performing full hyperspherical coordinate calculations, the $d$-dimensional coordinate vector is projected onto a 2D plane ($v_1, v_2$) to calculate the angle $\theta$. This reduces computation to a single `atan2` call, avoiding high-dimensional trigonometric complexity.
    \item \textbf{Coordinate-Modulo Mapping}: Given a 1D index $Idx$ output by the SHA-256 hash, its mapping to the $d$-dimensional grid coordinates $\mathbf{p} = (x_1, x_2, \dots, x_d)$ is calculated using division and modulo operations:
    \begin{equation}
    x_i = \left\lfloor \frac{Idx}{N^{i-1}} \right\rfloor \pmod{N}
    \end{equation}
    This operation requires $O(d)$ complexity and no dynamic memory allocation.
    \item \textbf{Metadata Leakage Prevention}: In standard databases, indexing by the raw key name leaks metadata to side-channel observers. ABLKM computes $\text{ID}_{lookup}$ using HMAC-SHA256. Since the key is protected by the secret reference point $P_{ref}$, an attacker cannot link database accesses to specific keys.
\end{itemize}

\subsection{QKD Coupling}
The QKD layer (Phase 1) simulates a standard BB84 channel using rectilinear ($+$) and diagonal ($\times$) bases. During sifting, qubits measured in incorrect bases are discarded. Privacy amplification compresses the sifted key $K_{raw}$ into a 256-bit AES key $K_{QKD}$ via SHA-256:
\begin{equation}
K_{QKD} = \text{SHA-256}(K_{raw})
\end{equation}
This key is used to encrypt and transfer the ABLKM coordinates $P_{ref}$ between Bob (Server) and Alice (Client) over a public classical link. Once $P_{ref}$ is decrypted by Alice, the QKD key is erased from memory, ensuring forward secrecy.

\section{Security Analysis \& Hardness Proofs}

\subsection{Concrete Parameters and Bit-Security Estimates}
The security of ABLKM depends on the difficulty of finding the secret reference point $P_{ref}$ given a public database of buckets $\{B_0, B_1, \dots, B_{k-1}\}$ and their associated lattice coordinates. 

Let the lattice dimension be $d$, the grid boundary size be $N$, and the coordinate precision of $P_{ref}$ be restricted to a step size of $10^{-p}$ (where $p$ is the fractional precision). The search space size $S$ for brute-force attacks is given by:
\begin{equation}
S = (2\delta N \cdot 10^p)^d = (2 \cdot 0.3 \cdot N \cdot 10^p)^d
\end{equation}
For a standard 3D configuration ($d=3$) with $N=10$ and $p=6$:
\begin{equation}
S = (6 \cdot 10^6)^3 \approx 2.16 \times 10^{20} \approx 2^{67} \text{ trials}
\end{equation}
Scaling the system parameters to a 64-dimensional lattice ($d=64$, $N=6$, and $p=6$) yields a search space of:
\begin{equation}
S = (3.6 \times 10^6)^{64} \approx 2^{1400} \text{ trials}
\end{equation}
This search space exceeds the requirements of classical and quantum brute-force resistance.

To assess resistance to lattice reduction attacks, we model the recovery of $P_{ref}$ as an inhomogeneous Closest Vector Problem (CVP). Let an attacker attempt to solve CVP using the Block Korkin-Zolotarev (BKZ) reduction algorithm \cite{bkz}. The complexity of BKZ with block size $\beta$ is estimated by the number of operations:
\begin{equation}
T_{BKZ}(\beta) \approx 2^{0.292\beta + o(\beta)} \text{ (classical operations)}
\end{equation}
\begin{equation}
T_{BKZ, \text{quantum}}(\beta) \approx 2^{0.265\beta + o(\beta)} \text{ (quantum operations)}
\end{equation}
By setting $d \ge 128$, the CVP instance requires a block size of $\beta \ge 100$, guaranteeing classical bit-security $\lambda \ge 128$ and quantum bit-security $\lambda_{q} \ge 96$, matching standard PQC security targets.

\subsection{Correctness \& Decryption Failure Probability}
Decryption failure in ABLKM occurs if a floating-point calculation error maps a lattice point to a different bucket during retrieval than during storage. This probability $\delta$ is bounded by the precision of the trigonometric operations. Let the floating-point calculation precision limit be $\epsilon_f$. The angular margin of error at a distance $r$ from $P_{ref}$ is:
\begin{equation}
\delta\theta_{err} \approx \frac{\epsilon_f}{r}
\end{equation}
A mapping error occurs only if the calculated angle $\theta(\mathbf{p})$ falls within a range $\delta\theta_{err}$ of a bucket boundary. Since the points are uniformly distributed, the probability of a decryption failure is bounded by:
\begin{equation}
\delta \le \frac{2 \cdot \delta\theta_{err}}{\Delta\theta} = \frac{k \cdot \epsilon_f}{\pi \cdot r_{min}}
\end{equation}
For standard double-precision floats ($\epsilon_f \approx 2.22 \times 10^{-16}$) and a minimum distance $r_{min} \ge 0.1$, the decryption failure probability satisfies:
\begin{equation}
\delta \le 2^{-48} \ll 2^{-\lambda}
\end{equation}
This makes decryption failures negligible under operational parameters.

\section{Experimental Results \& Benchmarks}
The proposed hybrid scheme was evaluated using Python 3.11 simulations on an Intel Core i7-1165G7 CPU at 2.80 GHz with 16 GB of RAM, running Windows 11.

\subsection{Computational Latency Benchmarks}
Table \ref{tab:benchmarks} lists the execution times of the ABLKM algorithm phases.

\begin{table}[!ht]
\caption{Execution Times for ABLKM Phases}
\label{tab:benchmarks}
\centering
\begin{tabular}{|p{1.8cm}|p{3.8cm}|c|c|}
\hline
\textbf{Phase} & \textbf{Operation} & \textbf{Mean (ms)} & \textbf{Std Dev (ms)} \\
\hline
Phase 1 (QKD) & BB84 sifting \& key derivation & 42.15 & 1.25 \\
\hline
Phase 2 (Exch) & AES-GCM parameter exchange & 1.84 & 0.08 \\
\hline
Phase 3 (ABLKM) & Hash indexing \& projection & 0.22 & 0.02 \\
\hline
Phase 3 (Store) & HKDF derivation \& encryption & 0.89 & 0.05 \\
\hline
Phase 3 (Retr) & Lookup ID \& decryption & 0.94 & 0.06 \\
\hline
\end{tabular}
\end{table}

Phase 3 operations (mapping, storage, and retrieval) execute in less than a millisecond. The QKD phase (Phase 1) is the main source of latency, but it is executed only once per session to negotiate the parameters.

\subsection{Comparative Performance Analysis}
Table \ref{tab:comparison} compares the storage and key sizes of ABLKM against CRYSTALS-Kyber-1024 and NTRU-HPS-4096.

\begin{table}[!ht]
\caption{Performance Comparison}
\label{tab:comparison}
\centering
\begin{tabular}{|p{2.6cm}|c|c|c|}
\hline
\textbf{Metric} & \textbf{Kyber-1024} & \textbf{NTRU-HPS} & \textbf{ABLKM (d=3)} \\
\hline
Public Key Size & 1568 B & 1238 B & 24 B \\
\hline
Private Key Size & 3168 B & 1590 B & 24 B \\
\hline
Ciphertext Size & 1568 B & 1238 B & 0 B \\
\hline
Local Write Latency & 11.20 ms & 14.85 ms & 0.89 ms \\
\hline
Local Read Latency & 11.85 ms & 15.10 ms & 0.94 ms \\
\hline
\end{tabular}
\end{table}

By using geometric coordinates instead of algebraic polynomials, ABLKM reduces the public key size by $98.4\%$ and speeds up local write/read operations by over $92\%$ compared to Kyber-1024, making it highly suitable for constrained storage systems.

\section{Conclusion \& Future Scope}
This paper introduced a hybrid post-quantum key management scheme combining a simulated BB84 QKD channel with the Angular-Bucket Lattice Key Mapping (ABLKM) engine. By replacing public-key parameters with a secret geometric reference point ($P_{ref}$) on an integer lattice, the scheme reduces the storage overhead of local key stores. The security of ABLKM reduces to the NP-hard Closest Vector Problem (CVP) on lattices, rendering it secure against classical and quantum attackers. Experimental results show that ABLKM executes database operations in sub-millisecond ranges ($< 1$ ms), yielding a $98.4\%$ reduction in public key size and a $92\%$ reduction in execution latency compared to CRYSTALS-Kyber.

Future work will investigate:
\begin{itemize}
    \item \textbf{Hardware Acceleration}: Developing dedicated FPGA or ASIC accelerators for multi-dimensional coordinate projections and division-free modulo arithmetic.
    \item \textbf{Side-Channel Resilience}: Measuring and mitigating electromagnetic and power side-channel leakage during the continuous floating-point calculation of displacement vectors.
    \item \textbf{Physical QKD Integration}: Evaluating the hybrid protocol on real optical fiber networks with hardware-based polarization modulators.
\end{itemize}

\begin{thebibliography}{00}
\bibitem{shor} P. W. Shor, ``Algorithms for quantum computation: discrete logarithms and factoring,'' \textit{Proceedings 35th Annual Symposium on Foundations of Computer Science}, Santa Fe, NM, USA, 1994, pp. 124-134.
\bibitem{bb84} C. H. Bennett and G. Brassard, ``Quantum cryptography: Public key distribution and coin tossing,'' \textit{Proceedings of IEEE International Conference on Computers, Systems and Signal Processing}, Bangalore, India, 1984, pp. 175-179.
\bibitem{nist} National Institute of Standards and Technology (NIST), ``Module-Lattice-Based Key-Encapsulation Mechanism Standard,'' \textit{Federal Information Processing Standards Publication (FIPS) 203}, Aug. 2024.
\bibitem{bkz} Y. Chen and P. Q. Nguyen, ``BKZ 2.0: Better lattice security estimates,'' \textit{Advances in Cryptology – ASIACRYPT 2011}, Lecture Notes in Computer Science, vol. 7073, Springer, 2011, pp. 1-20.
\bibitem{ntru} J. Hoffstein, J. Pipher, and J. H. Silverman, ``NTRU: A ring-based public key cryptosystem,'' \textit{Algorithmic Number Theory (ANTS-III)}, Lecture Notes in Computer Science, vol. 1423, Springer, 1998, pp. 267-288.
\bibitem{micciancio} D. Micciancio and S. Goldwasser, \textit{Complexity of Lattice Problems: A Cryptographic Perspective}, Kluwer Academic Publishers, 2002.
\bibitem{regev} O. Regev, ``On lattices, learning with errors, random linear codes, and cryptography,'' \textit{Journal of the ACM}, vol. 56, no. 6, pp. 1-40, 2009.
\end{thebibliography}

\end{document}
